% Quantitative receivers for the selected trace chart.
\section{Weighted price control for the trace receiver}
\label{sec:ex-lattice-trace-weighted-estimates}

This section proves the scalar combination of the response estimates that
feeds the rank-one receiver lemma
\ref{lem:ex-lattice-trace-weighted-price}. It isolates the field decay,
unweighted price convergence, and weighted price convergence that the
finite-volume producer must furnish. It does not assume the receiver itself
or suppress those three estimates into a single unnamed bound.

\begin{lemma}[Weighted price from a rank-one approximate kernel]
\label{lem:ex-lattice-trace-weighted-price}
Let $E$ and $H$ be finite-dimensional real normed spaces. Let
$f:E\to H$ and $U:H\to E$ be linear, with $fU=\mathrm{id}_H$ and
$\ker f=\mathbb R x$ for some $x\ne0$. Suppose linear maps
$v_i:E\to H$ and $u_i:H\to E$ satisfy
\[
 v_i\longrightarrow f,\qquad u_i\longrightarrow U,
 \qquad v_i u_i=\mathrm{id}_H\quad\text{eventually}.
\]
Let $w_i\to x$, let $e_i\in\mathbb R$, and let $\rho:E\to\mathbb R$
be a fixed linear functional. If
\[
 e_i v_i(w_i)\longrightarrow0\quad\text{in }H,
 \qquad e_i\rho(w_i)\longrightarrow0,
\]
then there are linear functionals $\mu_i:H\to\mathbb R$ with uniformly
bounded operator norms and
\[
 \bigl\|e_i(\mu_i v_i-\rho)\bigr\|_{E^*}\longrightarrow0.
\]
\end{lemma}
\begin{proof}
Choose a linear functional $\lambda:E\to\mathbb R$ with $\lambda(x)=1$,
by extending $x$ to a basis. Define
\[
 A(y)=\lambda\bigl(y-Uf(y)\bigr).
\]
Since $f(y-Uf(y))=0$, the kernel assumption gives
$y-Uf(y)=a x$ for some scalar $a$. Applying $\lambda$ shows $a=A(y)$,
and therefore
\[
 y=Uf(y)+A(y)x. \tag{*}
\]
Set $V_i=u_i f+A\otimes w_i$, where
$(A\otimes w_i)(y)=A(y)w_i$. The convergence of $u_i$ and $w_i$ implies
$V_i-\mathrm{id}_E=(u_i-U)f+A\otimes(w_i-x)$, by (*). The right side
tends to zero in operator norm. Once $\|V_i-\mathrm{id}_E\|<1$, the
Neumann series gives an inverse and
\[
 \|V_i^{-1}-\mathrm{id}_E\|
 \leq \frac{\|V_i-\mathrm{id}_E\|}{1-\|V_i-\mathrm{id}_E\|},
\]
so $V_i^{-1}\to\mathrm{id}_E$.

Put $\mu_i=\rho u_i$. These functionals have bounded norms because
$u_i\to U$. For $y\in E$, the exact identity $v_i u_i=\mathrm{id}_H$
gives
\begin{align*}
 e_i(\mu_i v_i-\rho)(V_i y)
 &=e_i\rho\bigl(u_i v_i(u_i f(y)+A(y)w_i)-u_i f(y)-A(y)w_i\bigr)\\
 &=A(y)\,e_i\bigl(\rho(u_i v_i(w_i))-\rho(w_i)\bigr).
\end{align*}
The first term inside the final difference vanishes after weighting,
because
\[
 e_i\rho(u_i v_i(w_i))=\rho(u_i(e_i v_i(w_i)))
\]
and $\rho u_i$ is uniformly bounded. The second term vanishes by hypothesis. Thus the
left side tends to zero for every fixed $y$. In finite dimension,
convergence on a basis implies operator-norm convergence, so
$e_i(\mu_i v_i-\rho)V_i\to0$. Multiplication by
$V_i^{-1}\to\mathrm{id}_E$ gives
$e_i(\mu_i v_i-\rho)\to0$ in operator norm.
\end{proof}

Let $E$ be a finite-dimensional real normed space, $H$ a finite-dimensional
real normed space, $s_i:E\to H$ linear maps, and $\rho:E\to\mathbb R$ a
fixed linear functional. Let $x\in E$ satisfy $\rho(x)=0$. For each $i$,
let $b_i\ge1$ with $b_i\to\infty$, let $q_i,r_i\in E$, and let
$e_i\ge0$. Define
\[
 m_i=b_i^{-1}q_i-r_i.
\]
Assume
\begin{align}
 b_i^{-1}q_i&\longrightarrow0,
 &\sup_i|\rho(q_i)|&<\infty,
 &r_i&\longrightarrow -x,
 \label{eq:ex-lattice-trace-response-limit}\\
 e_i/b_i&\longrightarrow0,
 &e_i\rho(r_i+x)&\longrightarrow0,
 &b_i s_i(m_i)&\longrightarrow0\quad\text{in }H.
 \label{eq:ex-lattice-trace-response-weighted}
\end{align}

\begin{lemma}[The medium has both weighted null limits]
\label{lem:ex-lattice-trace-medium-receipts}
Under \eqref{eq:ex-lattice-trace-response-limit}--
\eqref{eq:ex-lattice-trace-response-weighted},
\[
 m_i\longrightarrow x,\qquad
 e_i\rho(m_i)\longrightarrow0,\qquad
 e_i s_i(m_i)\longrightarrow0.
\]
\end{lemma}
\begin{proof}
Since $b_i^{-1}q_i\to0$ and $r_i\to-x$,
$m_i=b_i^{-1}q_i-r_i\to x$. Since $\rho(x)=0$,
\[
 e_i\rho(m_i)=(e_i/b_i)\rho(q_i)-e_i\rho(r_i+x).
\]
The first term tends to zero by boundedness of $\rho(q_i)$ and
$e_i/b_i\to0$, and the second by hypothesis. Also $e_i/b_i\le1$
eventually, so
\[
 \|e_i s_i(m_i)\|=(e_i/b_i)\|b_i s_i(m_i)\|
 \le\|b_i s_i(m_i)\|\longrightarrow0.
\]
\end{proof}

All sequences in the source application below are already reindexed by a
selected level $k_i\to\infty$: $b_i$ is the
arch scale at level $k_i$, $q_i$ and $r_i$ are the chosen and reference
responses there, and $s_i$ is the chart at index $i$. This reindexing matters:
the source does not assert the field-decay estimate at every unreindexed
level. Write $p_i=\widehat\beta_{k_i}$ and
$A_a(i)=\exp(a(\log\log p_i)^2)$. The source-specific weight is
$e_i=A_{1/5}(i)$ at this selected level. The next subsection proves
$e_i/b_i\to0$ and
$e_i\le p_i$. The field input is supplied with weight $p_i$; since
$b_i\le2p_i$, this gives $b_i s_i(m_i)\to0$ as required above.

In the selected construction, $E$ is the real signed-equivariant space of
grade-five coefficient arrays, $x$ is its nonzero equation-of-motion
direction, and $\rho$ is the free scalar price. The vectors $q_i$ and
$r_i$ are respectively the transported chosen response and the transported
reference response. Their combination $m_i=b_i^{-1}q_i-r_i$ is the
medium used to build a response functional. The scale $b_i$ is the positive
arch scale along the selected chart. Boundedness of $\rho(q_i)$ follows from the unweighted Ward-price limit;
the weighted estimate on $\rho(r_i+x)$ is the scale-resolved
reference-price comparison; and the last estimate is the field decay of
the medium.

Combining Lemma~\ref{lem:ex-lattice-trace-medium-receipts} with
Lemma~\ref{lem:ex-lattice-trace-weighted-price}, under the chart hypotheses
$s_i\to f$, $u_i\to U$, $fU=\mathrm{id}$, $\ker f=\mathbb Rx$ with
$x\ne0$, and $s_i u_i=\mathrm{id}$ eventually, gives bounded linear
functionals $\mu_i:H\to\mathbb R$ and
\[
 \sup_i\|\mu_i\|<\infty,
 \qquad
 \bigl\|e_i(\mu_i s_i-\rho)\bigr\|_{E^*}\longrightarrow0.
\]
This is the response-price estimate used by the two-direction trace
correction. The displayed deduction is complete; its chart and response
estimates are the separate producer obligations listed in the audit.

\subsection{The source rate and medium calculation}
\label{subsec:ex-lattice-trace-source-rate}

For the already selected levels $k_i\to\infty$, write
$p_i=\widehat\beta_{k_i}$ for the running positive rate and $b_i$ for the
arch scale at $k_i$. The scale comparison used by the construction is
\begin{equation}
 \frac12p_i\le b_i\le 2p_i,\qquad p_i\longrightarrow\infty.
 \label{eq:ex-lattice-trace-arch-comparison}
\end{equation}
The age weights are $A_a(i)=\exp(a(\log\log p_i)^2)$. They satisfy
\begin{equation}
 A_a(i)A_c(i)=A_{a+c}(i),\qquad
 A_a(i)\le p_i^\delta\quad\text{eventually for every fixed }a\in\mathbb R
 \text{ and every }\delta>0.
 \label{eq:ex-lattice-trace-age-rules}
\end{equation}
For the second relation, take logarithms. Since
$(\log\log p)^2=o(\log p)$, eventually
$|a|(\log\log p_i)^2\le\delta\log p_i$. Exponentiation gives the bound.
Also $A_a(i)\to0$ directly for $a<0$, since
$(\log\log p_i)^2\to\infty$.

In particular, for $a\in\{1/6,1/5\}$,
\begin{equation}
 0\le \frac{A_a(i)}{b_i}
 \le 2p_i^{a-1}\longrightarrow0.
 \label{eq:ex-lattice-trace-age-over-arch}
\end{equation}
The same argument gives $A_a(i)\le p_i$ eventually, so an estimate with
weight $p_i$ implies the corresponding estimate with weight $A_{1/5}(i)$.

Here is the scalar deduction from the response and reference receipts. Let
$q_i$ be the transported chosen response and $w_i$ the transported reference response. The local response estimate and the
separated flux limit have the form
\begin{equation}
 \|q_i\|\le C A_{1/6}(i),\qquad
 \rho(q_i)\longrightarrow-\kappa,
 \label{eq:ex-lattice-trace-response-inputs-source}
\end{equation}
while the reference chart and the reference-price comparison give
\begin{equation}
 w_i\longrightarrow -x,\qquad \rho(x)=0,\qquad
 A_{1/5}(i)\rho(w_i)\longrightarrow0.
 \label{eq:ex-lattice-trace-reference-inputs-source}
\end{equation}
Define $m_i=b_i^{-1}q_i-w_i$. Then
\[
 \|b_i^{-1}q_i\|\le 2C p_i^{-5/6}\longrightarrow0,
\]
so $m_i\to x$. Since $\rho(q_i)$ converges it is bounded, and
\[
 A_{1/5}(i)\rho(m_i)
 =\frac{A_{1/5}(i)}{b_i}\rho(q_i)
   -A_{1/5}(i)\rho(w_i)\longrightarrow0
\]
by \eqref{eq:ex-lattice-trace-age-over-arch} and
\eqref{eq:ex-lattice-trace-reference-inputs-source}. Finally, if the local
field estimate is $p_i\,s_i(m_i)\to0$, then
$A_{1/5}(i)\le p_i$ gives $A_{1/5}(i)s_i(m_i)\to0$. These are exactly
the three limits consumed by the rank-one approximate-kernel lemma; the
remaining price-receiver construction is its finite-dimensional argument,
not another analytic estimate.

The source establishes the arch comparison from positivity of the arch
scale and the two-sided scale control. It obtains the local response
estimate from a priced local charge inequality, transfers it through the
finite-volume separation limit, and gets the reference-price estimate from
the calibrated reference chart. Those three analytic inputs are not proved
by the scalar calculation above. Their declarations and proof boundaries are
listed in the adjacent evidence record.

\subsection{How the two weighted receipts are assembled}
\label{subsec:ex-lattice-trace-weighted-assembly}

The weighted price receipt for the medium uses only the response-price
limit, the reference-price cancellation, and the age-rate comparison just
proved. It is therefore useful to distinguish it from the stronger source
estimate that supplies the weighted field image. If $\rho(q_i)\to-\kappa$,
then $|\rho(q_i)|\le M$ eventually for some $M$; thus
\[
 \left|\frac{A_{1/5}(i)}{b_i}\rho(q_i)\right|
 \le 2M p_i^{-4/5}\longrightarrow0.
\]
No weighted convergence rate for the response flux itself is needed here.
The separate reference-chart estimate is the cancellation
\[
                         A_{1/5}(i)\rho(w_i)\to0.
\] This identifies the precise roles of the two
source limits and avoids conflating a field-size estimate with a price
estimate.

For the field image, the construction proves the stronger statement
$p_i s_i(m_i)\to0$ in the finite-dimensional chart norm. The transfer to
the weaker age weight is the monotonicity estimate
$A_{1/5}(i)\le p_i$; there is no loss from summing over a basis because the
chart is a fixed finite-dimensional operator. The bound
$\|q_i\|=O(A_{1/6}(i))$ alone proves only $m_i\to x$ and does not imply this
field estimate. The latter is supplied by the local pricing argument and
its finite-volume passage.

\subsection{Source route and exact remaining estimates}
\label{subsec:ex-lattice-trace-weighted-route}

The response-size and response-flux limits above follow from a concrete
separation argument. Let $\mathcal T$ be the linear map from the transported
finite-volume coefficient field to the grade-five response, and let
$\lambda$ be the flux functional on the transported fields. On the limiting
field, their compatibility is $\rho(\mathcal T v_i)=\lambda(v_i)$. At fixed
level $i$, write $v_i(p)$ for the finite-volume transported field indexed by the
configuration parameter $p$. The local charged-configuration estimate
gives, on a set in the separation ultrafilter,
\begin{align}
 \|\mathcal T v_i(p)\|&\le C_TA_{1/6}(i),\\
 |\lambda(v_i(p))+\kappa|&\le
 C_1A_{-1/4}(i)+C_2A_{-1}(i)
       +\frac{p_i^{-30}+p_i^{-100}}{12}.
 \label{eq:ex-lattice-trace-finite-response-price}
\end{align}
Here $C_T,C_1,C_2$ are fixed constants from the finite-dimensional
coefficient, response, and flux maps. The first line is the field-cost part
of the estimate; the second is the flux-price part.
The transported field converges along that ultrafilter to the infinite
volume field $v_i$. Since $\mathcal T$ and $\lambda$ are continuous linear
maps, the inequalities pass to their limits. Therefore
$q_i=\mathcal T v_i$ satisfies the response norm bound and
$\rho(q_i)=\lambda(v_i)\to-\kappa$. Every term on the right tends to zero by
\eqref{eq:ex-lattice-trace-age-rules}, proving
\eqref{eq:ex-lattice-trace-response-inputs-source}.

In the source, the finite-volume inequalities in
\eqref{eq:ex-lattice-trace-finite-response-price} are obtained by applying
the local charged-configuration bound to each typical configuration and
using the transported-field error estimate. The latter is bounded by
$A_{-1/4}$; the local flux estimate contributes
$4A_{-1}$ and the two inverse-power costs. Thus the ultrafilter passage is
an ordinary continuity argument; all quantitative work lies in the local
charge bound and its field-transport error estimate. Their exact source
declarations are recorded in the evidence table.

For the weighted reference-price estimate, let $c_i(p)$ be the calibrated
finite-volume probe and $v_i(p)$ the transported field. The calibration
gives $\rho(\mathcal T c_i(p))=\rho(-x)=0$, while the transport estimate
gives $\|v_i(p)-c_i(p)\|\le C A_{-1/4}(i)$ on an ultrafilter-large set.
Continuity therefore implies
\[
 |\rho(\mathcal T v_i(p))|
 \le \|\rho\|\,\|\mathcal T\| C A_{-1/4}(i).
\]
Passing to the ultrafilter limit gives the same bound for $|\rho(w_i)|$.
Multiplication by $A_{1/5}$ gives a quantity tending to zero because
$A_{1/5}A_{-1/4}=A_{-1/20}\to0$. This uses only a rate for the difference
between the transported field and the calibrated probe; it does not assert
that the probe itself converges at that rate.

The remaining chart-field input required by the general lemma is the exact
limit
\begin{equation}
 p_i\,s_i(m_i)\longrightarrow0.
 \label{eq:ex-lattice-trace-chart-field-decay}
\end{equation}
It is independent of the response-size and flux estimates above. To produce
it, the source constructs finite-volume labels whose one-point expectations
converge to the target core field, obtains local price bounds for each fixed
grade-four test, and uses derivative and orientation compatibility to pass
the bounds to a finite basis of the chart target. Finite dimensionality then
converts those basis estimates to the displayed chart-norm limit. The
finite-volume label existence, local price bounds, chart scaling identity,
and compatibility identities are analytic inputs to this step; no claim
that they follow from the scalar rate calculation is made here. Their exact
declarations and the source's application route are recorded only in the
evidence file.
